%% P11 --- Rozklad SVD, aproksymacja niskiego rzedu i uwarunkowanie
%%
%% Blizniak: programs/en/P11-svd-conditioning.tex. Ramka w ramke, makro w
%% makro, liczba w liczbe; tools/parity.py porownuje oba pliki.
%%
%% TEZA: kazda macierz, kwadratowa czy nie, rozklada sie na obrot-rozciagniecie
%% -obrot; obciecie rozciagniecia daje najlepsza aproksymacje danego rzedu; a
%% stosunek najwiekszego rozciagniecia do najmniejszego mowi, ile bledu solwer
%% wzmocni.
%%
%% P10 PRZEKAZUJE TEN PROGRAM PO IMIENIU, a jego trzy porazki sa tu motywacja.
%% P09 przekazuje polowe "nie odwracaj" dotyczaca dokladnosci -- splacone w
%% sekcji 5. P08 nazywa dlug w ramce ostrzezenia: rank collapse "potrzebuje
%% maszynerii z Programu P11 i prawdziwego modelu". Maszyneria jest tutaj.
%% Modelu nie ma, i sekcja 3 tak mowi, zamiast go wytwarzac.

\program{Rozkład SVD, aproksymacja niskiego rzędu i uwarunkowanie}
\label{prog:P11}
\index{singular value decomposition}
\index{matrix!condition number}

\begin{outcomes}
\outcome{1--6}{powiedzieć, co robią trzy czynniki rozkładu według wartości
  osobliwych i dlaczego każda macierz taki rozkład ma}
\outcome{7--11}{powiedzieć, dlaczego wartość osobliwa nie jest wartością
  własną i co ma zamiast niej macierz prostokątna}
\outcome{12--19}{powiedzieć, co zachowuje obcięcie do rzędu $r$ i ile ono
  kosztuje, oraz jak musi wyglądać widmo, żeby było tanie}
\outcome{20--24}{czytać wskaźnik uwarunkowania jako czynnik wzmocnienia błędu z
  wejścia na wyjście}
\outcome{25--33}{powiedzieć, dlaczego tworzenie $A\T A$ jest złym sposobem
  rozwiązania zadania najmniejszych kwadratów, wraz z cyframi}
\end{outcomes}

\begin{quiz}
\item Macierz przenosi wektor z przestrzeni jednego rozmiaru do przestrzeni
  innego. Czy może mieć wektor własny? \teachesat{1--2}
  \answerto{Nie \dash{} $Av = \lambda v$ wymaga, by $Av$ i $v$ żyły w tej samej
  przestrzeni, więc tego równania nie da się nawet zapisać. Wartości osobliwe
  są określone dla każdej macierzy każdego kształtu.}
\item $A$ ma wartości osobliwe $\val{p11.sig.hi}$ i $\val{p11.sig.lo}$. Podaj
  wartości własne $A\T A$. \teachesat{7--8}
  \answerto{$100$ i $4$ \dash{} kwadraty. $A\T A = V\Sigma^{2}V\T$, więc jej
  wartości własne to kwadraty wartości osobliwych, jakiekolwiek było $A$.}
\item Aproksymację rzędu $1$ odjęto od $A$. Co zostaje, w terminach wartości
  osobliwych? \teachesat{14--15}
  \answerto{Dokładnie $\sigma_2$, jeśli usunięty kawałek był największy. Każdy
  kawałek rzędu jeden wnosi własną wartość osobliwą i nic poza tym nie wnosi.}
\item $A$ ma wartości osobliwe między $\val{p11.sig.lo}$ a
  $\val{p11.sig.hi}$. O ile może wzmocnić błąd względny w najgorszym
  przypadku? \teachesat{20--21}
  \answerto{$\val{p11.kappa}$ razy \dash{} to stosunek największego
  rozciągnięcia do najmniejszego. Ten stosunek jest wskaźnikiem
  uwarunkowania.}
\item Tworzysz $A\T A$, żeby rozwiązać zadanie najmniejszych kwadratów. Co
  zrobiłeś ze wskaźnikiem uwarunkowania? \teachesat{25--26}
  \answerto{Podniosłeś go do kwadratu. $A\T A$ ma kwadraty wartości osobliwych,
  więc jej wskaźnik wynosi $\val{p11.kappa.sq}$ tam, gdzie dla $A$ wynosił
  $\val{p11.kappa}$ \dash{} a to kosztuje cyfry.}
\end{quiz}

% ============================================================
\section{Co przeżywa, gdy twierdzenie nie przeżywa}
\label{sec:P11-svd}
\index{singular value decomposition!definition}

\begin{fr}
Program~\ref{prog:P10} skończył się trzema macierzami, z którymi nie umiał
sobie poradzić, a nie były to macierze egzotyczne.

Macierz niesymetryczna, której wektory własne spotykały się pod kątem
$\val{p10.nonsym.angle}$ stopni zamiast prostego. Ćwierć obrotu bez żadnej
rzeczywistej wartości własnej. Ścinanie z jedną wartością własną dwa razy i
tylko jednym kierunkiem na pokaz.

I czwarty przypadek, którego on w ogóle nie podniósł, a który inżynier spotyka
najczęściej. Zapisz równanie wektora własnego $Av = \lambda v$ dla macierzy o
większej liczbie kolumn niż wierszy i powiedz, co idzie nie tak.
\dotline
\nextframe
\end{fr}

\begin{fr}
\begin{ansblock}
$Av$ i $v$ żyją w przestrzeniach różnych rozmiarów, więc tego równania nie da
się w ogóle zapisać.
\end{ansblock}

Wektor własny to kierunek, który wraca wzdłuż samego siebie, a \enquote{wzdłuż
samego siebie} wymaga, by wejście i wyjście były w tym samym pokoju. Macierz
$\val{p11.rect.rows} \times \val{p11.rect.cols}$ przenosi wektory z przestrzeni
trójwymiarowej do dwuwymiarowej. Niczego, co produkuje, nie da się porównać z
tym, co pochłonęła.

Cały Program~\ref{prog:P10} stosuje się więc do macierzy kwadratowych, a
większość macierzy w sieci kwadratowa nie jest.
\end{fr}

\begin{fr}
Naprawa zachowuje obraz i rezygnuje z ograniczenia.

Program~\ref{prog:P10} żądał \emph{jednego} zestawu prostopadłych kierunków,
które macierz jedynie skaluje. To dużo: te same kierunki muszą działać na obu
końcach.

Zażądaj zamiast tego \textbf{dwóch} zestawów \dash{} prostopadłych kierunków w
przestrzeni wejściowej i prostopadłych kierunków w przestrzeni wyjściowej, przy
czym macierz przenosi $i$-ty z pierwszych na wielokrotność $i$-tego z drugich:
\[ Av_i = \sigma_i u_i \]

Każda macierz to ma, każdego kształtu, bez wyjątków i bez warunków. Liczby
$\sigma_i$ to jej \textbf{wartości osobliwe}.
\end{fr}

\begin{fr}
Zapisane jako rozkład, z wektorami $v$ jako kolumnami $V$ i wektorami $u$ jako
kolumnami $U$:
\[ A = U\Sigma V\T \]
gdzie $U$ i $V$ są ortogonalne \dash{} słowo z Programu~\ref{prog:P09}: ich
transpozycje są ich odwrotnościami \dash{} a $\Sigma$ jest diagonalna z
wartościami osobliwymi na przekątnej, umownie od największej.

Czytane od prawej do lewej to trzy ruchy. $V\T$ obraca przestrzeń wejściową
tak, by wyróżnione kierunki ułożyły się wzdłuż osi. $\Sigma$ rozciąga każdą oś
własną liczbą. $U$ obraca wynik w to, gdzie mieszka wyjście.

\textbf{Obrót, rozciągnięcie, obrót.} Każda macierz, bez wyjątku.
\end{fr}

\mermaidfig{p11-three-moves}{Jedyne trzy rzeczy, jakie robi każda macierz, w
kolejności, w jakiej je robi.}{SVD jako trzy ruchy}

\begin{fr}
Program~\ref{prog:P10} sprawdził swoje twierdzenie dokładnie, budując
odpowiedź zamiast jej szukać, i ta sama sztuczka działa tutaj z jedną zmianą
wartą zauważenia.

Tam $A = QDQ\T$ używało \emph{tego samego} $Q$ po obu stronach, bo dwa zestawy
kierunków macierzy symetrycznej są tym samym zestawem. Tutaj nie muszą być,
więc konstrukcja używa dwóch różnych obrotów wymiernych \dash{} jednego z
trójkąta $3$--$4$--$5$ i jednego z $5$--$12$--$13$ \dash{} i tworzy
$A = U\Sigma V\T$, gdzie $\Sigma$ niesie $\val{p11.sig.hi}$ i
$\val{p11.sig.lo}$.

Powstałe $A$ ma wyrazy wybitnie nieprzyjemne, wśród nich $\frac{246}{65}$, i nie
jest symetryczne. Jego wartości osobliwe to dokładnie $\val{p11.sig.hi}$ i
$\val{p11.sig.lo}$, bo tam je wstawiono, a skrypt sprawdza
$Av_i = \sigma_i u_i$ dla obu $i$ na ułamkach.

Zanim pójdziesz dalej, przewidź o nim jedną rzecz. Wiedząc tylko, że jego
wartości osobliwe to $\val{p11.sig.hi}$ i $\val{p11.sig.lo}$, ile wynosi
$\lvert\det A\rvert$?
\dotline
\nextframe
\end{fr}

\begin{fr}
\begin{ansblock}
$\val{p11.detA}$, czyli
$\val{p11.sig.hi} \times \val{p11.sig.lo}$.
\end{ansblock}

To nie przypadek i nie jest to nic nowego. $\det(U\Sigma V\T) = \det U \cdot
\det\Sigma \cdot \det V\T$ na mocy reguły iloczynu z Programu~\ref{prog:P09},
oba obroty mają wyznacznik $\pm 1$, a $\det\Sigma$ jest iloczynem wartości
osobliwych. \textbf{Wyznacznik jest iloczynem rozciągnięć, a obroty nie wnoszą
nic} \dash{} czyli to, co Program~\ref{prog:P09} nazywał objętością,
powiedziane w języku tego programu.
\end{fr}

% ============================================================
\section{Nie wartości własne}
\label{sec:P11-not-eigen}
\index{singular value decomposition!against eigenvalues}

\begin{fr}
Te dwa pojęcia stoją dość blisko siebie, by mylić je bez przerwy, więc oto
zależność, dokładnie.

Weź $A = U\Sigma V\T$ i utwórz $A\T A$. Wymnóż, używając $U\T U = I$.
\dotline
\nextframe
\end{fr}

\begin{fr}
\begin{ansblock}
$A\T A = V\Sigma\T U\T U \Sigma V\T = V\Sigma^{2}V\T$.
\end{ansblock}

Czyli dokładnie to, co Program~\ref{prog:P10} zapisał jako $QDQ\T$, z $V$ w
miejsce $Q$ i $\Sigma^{2}$ w miejsce $D$.

Zatem: \textbf{wartości osobliwe $A$ to pierwiastki z wartości własnych
$A\T A$}, a prawe wektory osobliwe to wektory własne tamtej macierzy. $A\T A$
jest symetryczna, jakiekolwiek było $A$, więc twierdzenie spektralne zawsze się
do niej stosuje, a każde $\sigma_i$ jest rzeczywiste i nieujemne, bo jest
pierwiastkiem z liczby nieujemnej.

Dlatego wartości osobliwe istnieją dla każdej macierzy: są pożyczone z macierzy
symetrycznej, którą zawsze da się zbudować.
\end{fr}

\begin{fr}
Teraz pułapka, a ten przykład dobrano tak, by była trudna, a nie łatwa.

Zbudowane $A$ ma wartości osobliwe $\val{p11.sig.hi}$ i $\val{p11.sig.lo}$.
Jest kwadratowe, więc ma też własne wartości własne.

Zapisz, jakie według ciebie są, i czytaj dalej.
\dotline
\nextframe
\end{fr}

\begin{fr}
\begin{ansblock}
$\val{p11.lam.hi}$ i $\val{p11.lam.lo}$. Nie $\val{p11.sig.hi}$ i
$\val{p11.sig.lo}$.
\end{ansblock}

Blisko. Na tyle blisko, że na wykresie różnicy byś nie zobaczył, i na tyle
blisko, że tabela podająca niewłaściwą parę przeszłaby przez recenzję.

\transcript{p11-singular-vs-eigen}

Listing idzie drugą drogą: ślad i wyznacznik $A\T A$ to $104$ i $400$, czyli
suma i iloczyn liczb $100$ i $4$ \dash{} kwadratów wartości osobliwych,
dokładnie.
\end{fr}

\begin{trapbox}
\textbf{\enquote{Wartości osobliwe to wartości własne.}}

Pokrywają się dla jednej rodziny: macierzy symetrycznych dodatnio
półokreślonych, gdzie $A\T A = A^{2}$ i pierwiastek cofa kwadrat. Do tej
rodziny należą macierze kowariancji i macierze Grama, i dlatego nawyk się
wyrabia.

Wszędzie indziej różnią się, i to nie tylko wartością, ale rodzajem:

\begin{itemize}
\item wartości osobliwe istnieją dla \textbf{każdej} macierzy, także
  prostokątnej; wartości własne wymagają kwadratowej;
\item wartości osobliwe są \textbf{rzeczywiste i nieujemne}, zawsze; wartości
  własne bywają zespolone, jak pokazała ćwierć obrotu z
  Programu~\ref{prog:P10};
\item wartości osobliwych jest zawsze \textbf{tyle, ile wynosi mniejszy
  wymiar}, z pełnym zestawem kierunków; wektorów własnych może zabraknąć, jak
  pokazało ścinanie.
\end{itemize}

Powyższy przykład jest przypadkiem niewygodnym: $\val{p11.lam.hi}$ wobec
$\val{p11.sig.hi}$ to różnica $5\%$, a wykres jednego opisany jako drugie
wyglądałby zupełnie rozsądnie.
\end{trapbox}

\begin{fr}
Dla macierzy prostokątnej są dwie macierze Grama zamiast jednej, a rozmiary
warto zobaczyć.

Macierz $\val{p11.rect.rows} \times \val{p11.rect.cols}$ o nazwie $R$ daje
$R\T R$ rozmiaru $\val{p11.rect.cols} \times \val{p11.rect.cols}$ oraz $RR\T$
rozmiaru $\val{p11.rect.rows} \times \val{p11.rect.rows}$. Różne macierze,
różne rozmiary, i żadna z nich nie jest $R$.

Ich ślady zgadzają się na $\val{p11.rect.energy}$, co skrypt sprawdza, i to
jest wskazówka: obie niosą tę samą sumę, więc implikowane przez nie wartości
osobliwe są tą samą listą. \textbf{Większa macierz Grama ma po prostu dodatkowe
zera}, po jednym na każdy wymiar, do którego macierz nie sięga.
\end{fr}

% ============================================================
\section{Zachowywanie największych kawałków}
\label{sec:P11-lowrank}
\index{matrix!low-rank approximation}

\begin{fr}
Rozkład można czytać na drugi sposób, i to ten drugi jest użyteczny.

$U\Sigma V\T$ wymnaża się do \emph{sumy}:
\[ A = \sigma_1 u_1v_1\T + \sigma_2u_2v_2\T + \cdots \]
Każdy składnik to kolumna razy wiersz, co na mocy Programu~\ref{prog:P08} jest
macierzą rzędu dokładnie jeden. Macierz jest więc sumą kawałków rzędu jeden, po
jednym na wartość osobliwą, z nią jako wagą.

Kawałki nie oddziałują: wektory $u$ są prostopadłe i wektory $v$ też. Powiedz,
co to daje \dash{} na co rozkłada się całkowity \enquote{rozmiar} $A$?
\dotline
\nextframe
\end{fr}

\begin{fr}
\begin{ansblock}
Na kwadraty wartości osobliwych i nic poza tym.
\end{ansblock}

Suma kwadratów wszystkich wyrazów $A$ \dash{} wielkość, którą
Program~\ref{prog:P05} nazwałby kwadratem długości $A$ czytanego jako jeden
długi wektor \dash{} wynosi $\sigma_1^{2} + \sigma_2^{2} + \cdots$.

Prostopadłe kawałki dodają się w kwadratach, co jest twierdzeniem Pitagorasa z
Programu~\ref{prog:F09} z macierzami w miejsce strzałek. Każda wartość osobliwa
niesie więc swój udział w całości, a udziałem jest $\sigma_i^{2}$.
\end{fr}

\mermaidfig{p11-keep-the-largest}{Dlaczego kawałki macierzy można ważyć po
jednym.}{macierz jako suma kawalkow}

\begin{fr}
Teraz aproksymacja. Zachowaj pierwszych $r$ składników, a resztę wyrzuć.

Właśnie ci powiedziano, co wnosi każdy kawałek. Zatem bez liczenia czegokolwiek:
co zostaje, mierzone tak samo?
\dotline
\nextframe
\end{fr}

\begin{fr}
\begin{ansblock}
$\sigma_{r+1}^{2} + \sigma_{r+2}^{2} + \cdots$ \dash{} dokładnie te kawałki,
które wyrzuciłeś.
\end{ansblock}

Na zbudowanej macierzy z sekcji 1 obcięcie do rzędu $1$ zostawia błąd dokładnie
$\val{p11.ey.err}$, czyli $\sigma_2$, i skrypt sprawdza to na ułamkach.

Błąd nie jest \emph{mniej więcej} $\sigma_2$ ani przez tę wielkość
ograniczony. Jest $\sigma_2$.

Co prowadzi do pytania, o które chodzi w reszcie sekcji. Obcięcie to
\emph{jeden} sposób zbudowania macierzy rzędu $1$ blisko $A$. Czy jest lepszy?
\dotline
\nextframe
\end{fr}

\begin{fr}
\begin{ansblock}
Nie. Nie ma macierzy rzędu $1$ bliższej $A$ niż tamta.
\end{ansblock}

\textbf{Twierdzenie Eckarta--Younga}, które ta książka podaje i którego nie
dowodzi. Ze wszystkich macierzy rzędu $r$ obcięty rozkład SVD jest najbliższy
$A$. Nie dobry wybór, nie dająca się obronić heurystyka: najlepszy, jaki
istnieje.

To mocna teza \dash{} kwantyfikuje po każdej macierzy rzędu $r$, a jest ich
nieskończenie wiele \dash{} więc zasługuje na to samo traktowanie, jakie
Program~\ref{prog:P10} dał twierdzeniu spektralnemu. Skrypt buduje
$\val{p11.ey.trials}$ innych macierzy rzędu jeden z wymiernych par kierunków,
bierze najlepsze skalowanie dla każdej i wymaga, by każda była co najmniej tak
samo zła. Każda jest.
\end{fr}

\begin{fr}
Jak w Programie~\ref{prog:P10}, te próby są dowodem o kodzie, a nie o
twierdzeniu: kontrprzykład obaliłby dowód.

Do czego służą, jest tu jednak inne i warto to oddzielić. Twierdzenie
spektralne sprawdzano, bo słowo \emph{dokładnie} wymagało dokładnej arytmetyki.
Eckarta--Younga sprawdza się, bo \textbf{tę tezę łatwo zapamiętać jako
słabszą, niż jest} \dash{} \emph{dobra aproksymacja} zamiast \emph{najlepsza}
\dash{} a zobaczenie tysięcy alternatyw, które przegrywają, ustala, którą z
dwóch czytelnik wyniesie.
\end{fr}

\begin{fr}
Oto, co to daje, i to jest powód, dla którego obchodzi to kogokolwiek spoza
wydziału matematyki.

Weź macierz, której wartości osobliwe biegną $100, 40, 16, 6, 2, 1$. Zachowaj
$\val{p11.spec.keep}$ największe z $\val{p11.spec.n}$, a cztery wyrzuć.

Jaki ułamek całości nadal masz?
\dotline
\nextframe
\end{fr}

\begin{fr}
\begin{ansblock}
$\val{p11.spec.pct}\%$, zostawiając $\val{p11.spec.rest.pct}\%$.
\end{ansblock}

Jedna trzecia kawałków niosąca prawie całą treść. To cały argument za
przechowywaniem $U_r\Sigma_r V_r\T$ zamiast $A$ i za adapterem niskiego rzędu z
Programu~\ref{prog:P08}: jeśli macierz jest prawie rzędu
$\val{p11.spec.keep}$, traktowanie jej jako rzędu $\val{p11.spec.keep}$
kosztuje prawie nic.

Zauważ kształt warunku. Nie chodzi o rozmiar macierzy. Chodzi o to, jak szybko
wartości osobliwe \textbf{zanikają}, a macierz o płaskim widmie nie zyskuje
nic.
\end{fr}

\begin{warning}
\textbf{Tamto widmo zostało zbudowane, a ta książka nie zmierzyła żadnego
prawdziwego.}

Powyższa ramka wybrała sześć liczb, które szybko zanikają, i podała, co z tego
wynika. Pokazuje \emph{mechanizm} i nie jest dowodem, że jakakolwiek konkretna
macierz tak się zachowuje.

Teza empiryczna \dash{} że macierze zanurzeń albo wag wytrenowanego modelu mają
widma skupione na tyle, by aproksymacja rzędu $r$ była niemal darmowa \dash{}
jest tym, na czym opiera się literatura o adaptacji niskiego rzędu, a
rozstrzygnięcie jej wymaga wartości osobliwych prawdziwych macierzy prawdziwego
modelu. Ta książka takiego nie ma i nie pobiera.

To ten sam dług, którego Program~\ref{prog:P08} nie udawał, że spłacił, w
sprawie rank collapse, i ta sama odpowiedź: maszyneria jest tutaj, modelu nie
ma, a powiedzenie, co jest czym, jest całą tą dyscypliną.
\end{warning}

% ============================================================
\section{Błąd na wejściu, błąd na wyjściu}
\label{sec:P11-conditioning}
\index{matrix!condition number}

\begin{fr}
Inne pytanie o te same trzy ruchy, i to ono rozstrzyga, czy odpowiedzi można
ufać.

Wejście przychodzi już obarczone błędem \dash{} z pomiaru, z czujnika albo po
prostu z ostatniego zaokrąglenia, które opisał Program~\ref{prog:P01}. Macierz
działa na błąd dokładnie tak, jak działa na wszystko inne.

Jeśli błąd wejścia akurat wskazuje wzdłuż $v_1$, zostaje pomnożony przez
$\sigma_1$. Jeśli wskazuje wzdłuż ostatniego $v$, przez najmniejszą $\sigma$.
Pytanie brzmi, co dzieje się z błędem \emph{względem} odpowiedzi, a to jest
stosunek dwóch z nich.
\dotline
\nextframe
\end{fr}

\begin{fr}
\begin{ansblock}
W najgorszym przypadku błąd względny zostaje pomnożony przez
$\sigma_{\max}/\sigma_{\min}$.
\end{ansblock}

W najgorszym przypadku odpowiedź jest tak mała, jak odwzorowanie pozwala, a
błąd tak duży, jak pozwala, więc oba skrajne przypadki się dzielą.

Ten stosunek to \textbf{wskaźnik uwarunkowania}, pisany $\kappa(A)$, a dla
macierzy z sekcji 1 wynosi
$\val{p11.sig.hi}/\val{p11.sig.lo} = \val{p11.kappa}$.

Jest bezwymiarowy i mnoży błąd \emph{względny}. Zatem: wejście dobre do
szesnastu cyfr znaczących wchodzi w macierz o $\kappa = 10^{6}$. Do mniej
więcej ilu cyfr dobra jest odpowiedź?
\dotline
\nextframe
\end{fr}

\mermaidfig{p11-error-in-error-out}{Co macierz robi z błędem, który przyszedł
razem z wejściem.}{wskaznik uwarunkowania jako wzmocnienie}

\begin{fr}
\begin{ansblock}
Do jakichś dziesięciu. Sześć z szesnastu przepadło.
\end{ansblock}

Reguła kciuka jest tą, którą warto nieść: \textbf{tracisz mniej więcej
$\logb{10}\kappa$ cyfr dziesiętnych}. Podwójna precyzja zaczyna z jakimiś
szesnastoma, więc wskaźnik $10^{16}$ nie zostawia nic, a wskaźnik
$\val{p11.spec.kappa}$ kosztuje dwie z szesnastu i jest niczym szczególnym.
\end{fr}

\begin{fr}
Dwie rzeczy, którymi on nie jest, a obie są częstymi odczytami.

Nie jest własnością \emph{algorytmu}. Wskaźnik uwarunkowania jest faktem o
macierzy i mówi, z czym musi się zmierzyć każda poprawna metoda. Źle
uwarunkowane zadanie rozwiązane doskonałym algorytmem nadal daje kiepską
odpowiedź, bo odpowiedź naprawdę jest tak czuła na wejście.

I nie jest zgłoszeniem błędu. Wąwóz z Programu~\ref{prog:P10} miał duży
wskaźnik z założenia; ma go też każde zadanie, w którym jeden kierunek liczy
się znacznie bardziej niż inny. Ta liczba mówi, ile precyzji zadanie
kosztuje, a nie że ktoś popełnił pomyłkę.

Oba odczyty biorą się z tego samego miejsca: liczba zapowiadająca zły wynik
bywa traktowana jak diagnoza. Bliżej jej do prognozy pogody. Mówi, jakie są
warunki i ile będą kosztować, a uczciwą reakcją jest zaplanowanie tej straty, a
nie szukanie pomyłki, która ją spowodowała.

Co warto połączyć z czymś już zmierzonym. Dla macierzy symetrycznej dodatnio
określonej wartości osobliwe \emph{są} wartościami własnymi. Co więc
Program~\ref{prog:P10} nazwał już $\kappa$, nie używając tego słowa?
\dotline
\nextframe
\end{fr}

\begin{fr}
\begin{ansblock}
Stosunek wartości własnych niecki \dash{} którego pierwiastek mówił, jak
wydłużona jest elipsa poziomicowa.
\end{ansblock}

To daje połączenie, wokół którego krążyły dwa programy. Wąwóz to niecka o
dużym $\kappa$; duże $\kappa$ to zadanie oddające cyfry; a
Program~\ref{prog:P20} pokaże, że to także zadanie kosztujące optymalizator
iteracje.

\textbf{Jedna liczba, trzy konsekwencje}, a wkładem tego programu jest ta
środkowa.
\end{fr}

% ============================================================
\section{Dlaczego kod nie tworzy \texorpdfstring{$A\T A$}{A transponowane A}}
\label{sec:P11-normal}
\index{least squares!normal equations}

\begin{fr}
Program~\ref{prog:P09} zmierzył, że rozwiązanie układu kosztuje jakąś piątą
tego, co jego odwracanie, a potem powiedział coś, czego nie mógł poprzeć:

\begin{quote}
Rachunek operacji jest mierzony tutaj. Dokładność należy do
Programu~\ref{prog:P11} i jest tą połową, która rozstrzyga spór.
\end{quote}

Ta sekcja to spłaca. Nośnikiem są najmniejsze kwadraty, bo tam pokusa jest
największa: Program~\ref{prog:P08} wyprowadził równania normalne
$A\T A\hat{x} = A\T b$ jako postać zamkniętą, a postać zamknięta prosi się o
wpisanie tak, jak stoi.

Wiesz, jakie są wartości osobliwe $A\T A$. Ile więc wynosi jej wskaźnik
uwarunkowania?
\dotline
\nextframe
\end{fr}

\begin{fr}
\begin{ansblock}
$\kappa(A)^{2}$. Wartości osobliwe są podniesione do kwadratu, więc ich
stosunek też.
\end{ansblock}

Dla macierzy z sekcji 1 to $\val{p11.kappa.sq}$ wobec $\val{p11.kappa}$, co
skrypt sprawdza dokładnie.

A teraz zestaw to z regułą kciuka sprzed dwóch ramek. Jeśli tracisz mniej
więcej $\logb{10}\kappa$ cyfr, a tworzenie $A\T A$ podnosi $\kappa$ do
kwadratu, to tworzenie $A\T A$ \textbf{podwaja liczbę traconych cyfr}. Nie
trochę gorzej: dwa razy tyle straty, za krok zrobiony tylko dlatego, że wzór
tak został zapisany.
\end{fr}

\begin{fr}
To był argument. Oto pomiar, i jest to jedyna teza w tym programie naprawdę o
liczbach zmiennoprzecinkowych, a nie o macierzach.

Prawdziwe zadanie najmniejszych kwadratów: dopasuj wielomian stopnia
$\val{p11.ls.deg}$ przez $\val{p11.ls.points}$ punktów rozłożonych na
$\intcc{0}{1}$. Współczynniki są wybrane, więc dokładna odpowiedź jest znana, a
prawa strona zbudowana na ułamkach, żeby nic nie przepadło przed rozwiązywaniem.

Potem ten sam układ rozwiązuje się dwa razy w zwykłej podwójnej precyzji
\dash{} raz przez $A\T A$, raz rozkładem, który jej nigdy nie tworzy.
\end{fr}

\begin{fr}
\begin{tabularx}{\linewidth}{@{}lX@{}}
\toprule
\textbf{droga} & \textbf{błąd względny współczynników} \\
\midrule
rozkład, $A\T A$ nigdy nie utworzone & lepszy niż $\val{p11.ls.qr.bound}$ \\
równania normalne & gorszy niż $\val{p11.ls.ne.floor}$ \\
\bottomrule
\end{tabularx}

\medskip
Co najmniej $\val{p11.ls.digits}$ cyfr dziesiętnych, oddanych za wpisanie wzoru
tak, jak stoi. Obie wielkości podano jako ograniczenia, a nie jako pomiary, a
powód jest z Programu~\ref{prog:P01}: reszta jest własnością maszyny, która ją
wyprodukowała, więc tym, co się tu zapisuje, jest różnica, która własnością
maszyny nie jest.
\end{fr}

\begin{fr}
Ta porażka ma kształt, który czyni ją groźną, a nie kształt, który czyni ją
oczywistą.

Nic się nie podnosi. Żadne ostrzeżenie się nie pojawia. Współczynniki, które
wracają, wyglądają jak współczynniki: właściwy rząd wielkości, właściwe znaki,
wiarygodne dla każdego, kto nie znał odpowiedzi wcześniej. Są po prostu błędne
od szóstej cyfry znaczącej wzwyż.

Program~\ref{prog:P07} nazwał ten kształt dokładnie \dash{} usterka łagodna,
póki się debuguje, a dotkliwa, gdy wszystko zaczyna działać, to najgorszy
kształt usterki. Tutaj odpowiednikiem jest: usterka, która wygląda jak
odpowiedź, jest gorsza niż taka, która wygląda jak awaria.
\end{fr}

\begin{fr}
Rada \enquote{nie twórz równań normalnych} to więc dwie tezy z dwoma
właścicielami, i obie mają liczby.

Program~\ref{prog:P09}: tworzenie odwrotności kosztuje jakieś
$\val{p09.cost.ratio}$ razy tyle arytmetyki co rozwiązywanie, zmierzone
liczeniem. Ten program: tworzenie $A\T A$ podnosi wskaźnik uwarunkowania do
kwadratu, więc podwaja tracone cyfry, zmierzone dwukrotnym rozwiązaniem tego
samego układu.

Ta druga rozstrzyga. Czynnik $\val{p09.cost.ratio}$ w szybkości to kompromis
inżynierski, który wchłonie szybsza maszyna. Cyfr nie odzyskuje się przy żadnej
szybkości.
\end{fr}

\begin{fr}
Jeszcze jeden obiekt, i jest to definicja, a nie technika.

Program~\ref{prog:P08} pokazał, że wysoki układ zwykle nie ma rozwiązania i że
najlepszą odpowiedzią jest rzut. SVD daje tej odpowiedzi postać zamkniętą:
obróć przez $U\T$, podziel przez każdą wartość osobliwą, obróć przez $V$.
Zapisane jako macierz nazywa się to \textbf{pseudoodwrotnością}, $A^{+}$.

Tam, gdzie $A$ jest odwracalne, zgadza się z $A^{-1}$. Tam, gdzie nie jest,
daje odpowiedź najmniejszych kwadratów, a gdy kilka odpowiedzi remisuje, daje
najkrótszą.

Zauważ, że wszystkie trzy kroki to rzeczy, które ten program już nazwał. Obrót
macierzą ortogonalną to bezstratna zmiana bazy z Programu~\ref{prog:P09};
dzielenie przez wartości osobliwe to jedyne miejsce, w którym cokolwiek się
skaluje; a obrót z powrotem to ten sam ruch jeszcze raz. Nic w tej definicji
nie jest nowe, i dlatego jest to definicja, a nie technika.

Na ten środkowy krok warto popatrzeć dłużej. Dzieli przez \emph{każdą} wartość
osobliwą, także najmniejszą. Powiedz, co to robi w kierunku, który macierz
prawie spłaszczyła.
\dotline
\nextframe
\end{fr}

\begin{fr}
\begin{ansblock}
Prawie eksploduje: dzielenie przez coś prawie zerowego.
\end{ansblock}

To znowu wskaźnik uwarunkowania, przybywający jako dzielenie, które da się
wskazać palcem.

Dlatego też obcinanie pomaga: odrzucenie najmniejszych wartości osobliwych
przed odwracaniem ogranicza wzmocnienie, kosztem nieodpowiadania w tych
kierunkach w ogóle. Czy ten kompromis warto zawrzeć, jest osądem o zadaniu, a
ta książka nie ma na to reguły.
\end{fr}

\begin{fr}
Część III zamyka się w tym miejscu, i zamyka się jednym rozkładem.

Program~\ref{prog:P04} pytał, czym jest przestrzeń, a Program~\ref{prog:P06}
uczynił macierz funkcją. Program~\ref{prog:P08} policzył, jak dużo przestrzeni
ją przeżywa, Program~\ref{prog:P09} zmierzył ile, a Program~\ref{prog:P10}
znalazł kierunki przechodzące nieobrócone \dash{} dla macierzy, które je mają.

$A = U\Sigma V\T$ ma to wszystko. Rząd to liczba niezerowych wartości
osobliwych; wyznacznik to ich iloczyn; jądro rozpinają te wektory $v$, których
$\sigma$ jest zerem; najlepsza aproksymacja niskiego rzędu to kilka pierwszych
składników; a wskaźnik uwarunkowania to pierwsza podzielona przez ostatnią.
\textbf{Pytania sześciu programów, odpowiedziane odczytaniem jednego
rozkładu}, i to jest najmocniejszy powód, by wiedzieć, że ten rozkład istnieje.
\end{fr}

% ============================================================
\begin{summarybox}
\sumitem{1--2}{\textbf{Równania wektora własnego nie da się w ogóle zapisać dla
  macierzy prostokątnej}, bo $Av$ i $v$ żyją w przestrzeniach różnych rozmiarów.
  Program~\ref{prog:P10} stosuje się więc do macierzy kwadratowych, a większość
  macierzy w sieci kwadratowa nie jest.}
\sumitem{3--4}{Naprawa żąda \textbf{dwóch} zestawów prostopadłych kierunków
  zamiast jednego \dash{} $Av_i = \sigma_i u_i$ \dash{} i ma je każda macierz
  każdego kształtu. Zapisane jako $A = U\Sigma V\T$ jest to \textbf{obrót,
  rozciągnięcie, obrót}, czytane od prawej do lewej, bez wyjątków i bez
  warunków.}
\sumitem{5--6}{Sprawdzone dokładnie na macierzy zbudowanej z \emph{dwóch}
  różnych obrotów wymiernych, czyli konstrukcją z Programu~\ref{prog:P10} ze
  zniesionym warunkiem $U = V$. A $\lvert\det A\rvert$ jest iloczynem wartości
  osobliwych, bo obroty nie wnoszą nic \dash{} czyli objętość z
  Programu~\ref{prog:P09} w języku tego programu.}
\sumitem{7--8}{$A\T A = V\Sigma^{2}V\T$, więc \textbf{wartości osobliwe są
  pierwiastkami z wartości własnych $A\T A$}. Dlatego istnieją dla każdej
  macierzy: są pożyczone z macierzy symetrycznej, którą zawsze da się zbudować,
  i twierdzenie spektralne zawsze się do niej stosuje.}
\sumitem{9--10}{\textbf{Wartości osobliwe to nie wartości własne}, a przykład
  dobrano tak, by pomyłka była zrozumiała: $\val{p11.lam.hi}$ i
  $\val{p11.lam.lo}$ wobec $\val{p11.sig.hi}$ i $\val{p11.sig.lo}$, różnica
  $5\%$, której żaden wykres nie pokaże. Różnią się też rodzajem \dash{}
  wartości osobliwe istnieją dla każdego kształtu, są zawsze rzeczywiste i
  nieujemne i nigdy nie brakuje im kierunków.}
\sumitem{11}{Macierz prostokątna ma \textbf{dwie} macierze Grama różnych
  rozmiarów, których ślady zgadzają się na $\val{p11.rect.energy}$. Większa
  niesie po prostu dodatkowe zera, po jednym na każdy wymiar, do którego
  macierz nie sięga.}
\sumitem{12--13}{$A$ jest \textbf{sumą kawałków rzędu jeden}, po jednym na
  wartość osobliwą, i one nie oddziałują \dash{} całkowity rozmiar rozkłada się
  więc na $\sigma_i^{2}$, co jest twierdzeniem Pitagorasa z
  Programu~\ref{prog:F09} z macierzami w miejsce strzałek.}
\sumitem{14--15}{Obcięcie do rzędu $r$ zostawia zatem \textbf{dokładnie}
  wyrzucone kawałki: $\val{p11.ey.err}$ dla przykładowej macierzy, czyli
  $\sigma_2$. Nie mniej więcej $\sigma_2$, nie ograniczone przez nią.}
\sumitem{16--17}{\textbf{Eckart--Young}: ze wszystkich macierzy rzędu $r$
  obcięcie jest najbliższe. Podane, nie dowodzone, i sprawdzone wobec
  $\val{p11.ey.trials}$ alternatyw \dash{} nie jako dowód twierdzenia, tylko
  dlatego, że \textbf{tę tezę łatwo zapamiętać jako słabszą, niż jest}, a
  patrzenie, jak tysiące alternatyw przegrywa, ustala, którą wersję czytelnik
  wyniesie.}
\sumitem{18--19}{Dwie z sześciu wartości osobliwych niosące $\val{p11.spec.pct}\%$
  to cały argument za przechowywaniem obcięcia i za adapterem niskiego rzędu. A
  warunek dotyczy \textbf{tego, jak szybko widmo zanika}, a nie rozmiaru
  macierzy: płaskie widmo nie zyskuje nic.}
\sumitem{20--21}{Błąd przychodzi z wejściem i jest rozciągany jak wszystko inne,
  więc w najgorszym przypadku błąd \emph{względny} zostaje pomnożony przez
  $\sigma_{\max}/\sigma_{\min}$. To \textbf{wskaźnik uwarunkowania}, a dla
  przykładowej macierzy wynosi $\val{p11.kappa}$.}
\sumitem{22--23}{\textbf{Tracisz mniej więcej $\logb{10}\kappa$ cyfr
  dziesiętnych} z szesnastu, od których zaczyna podwójna precyzja. Nie jest to
  własność algorytmu \dash{} doskonała metoda nadal daje kiepską odpowiedź na
  źle uwarunkowanym zadaniu \dash{} i nie jest to zgłoszenie błędu.}
\sumitem{24}{Dla macierzy symetrycznej dodatnio określonej $\kappa$ jest
  stosunkiem wartości własnych, który zmierzył Program~\ref{prog:P10}, więc
  \textbf{wąwóz, utrata cyfr i wolny optymalizator to jedna liczba czytana na
  trzy sposoby}. Ta środkowa należy do tego programu.}
\sumitem{25--26}{$A\T A$ ma \emph{kwadraty} wartości osobliwych, więc
  \textbf{$\kappa(A\T A) = \kappa(A)^{2}$} \dash{} $\val{p11.kappa.sq}$ wobec
  $\val{p11.kappa}$. Podniesienie $\kappa$ do kwadratu podwaja tracone cyfry,
  za krok zrobiony tylko dlatego, że postać zamknięta z
  Programu~\ref{prog:P08} tak została zapisana.}
\sumitem{27--29}{Zmierzone na dopasowaniu stopnia $\val{p11.ls.deg}$ przez
  $\val{p11.ls.points}$ punktów: rozkład schodzi poniżej
  $\val{p11.ls.qr.bound}$, równania normalne nie schodzą poniżej
  $\val{p11.ls.ne.floor}$, \textbf{co najmniej $\val{p11.ls.digits}$ cyfr
  oddanych}. Zapisane jako ograniczenia, bo reszta należy do maszyny. I nic się
  nie podnosi: błędne współczynniki wyglądają jak współczynniki, co jest gorsze
  niż awaria.}
\sumitem{30}{Rada \enquote{nie twórz równań normalnych} to więc dwie tezy z
  dwoma właścicielami. Program~\ref{prog:P09} zmierzył arytmetykę, ten program
  dokładność \dash{} a \textbf{dokładność rozstrzyga}, bo czynnik w szybkości
  wchłonie szybsza maszyna, a cyfr nie odzyskuje się przy żadnej szybkości.}
\sumitem{31--32}{\textbf{Pseudoodwrotność} $A^{+}$ to obróć, podziel przez każdą
  wartość osobliwą, obróć \dash{} zgodna z $A^{-1}$ tam, gdzie ta istnieje, i
  dająca odpowiedź najmniejszych kwadratów tam, gdzie nie. Dzielenie przez
  najmniejszą $\sigma$ jest miejscem, w którym wskaźnik uwarunkowania staje się
  dzieleniem, które da się wskazać palcem, a wcześniejsze obcięcie jest
  kompromisem, który je ogranicza.}
\sumitem{33}{\textbf{Jeden rozkład odpowiada na pytania sześciu programów}: rząd
  to liczba niezerowych $\sigma$, wyznacznik to ich iloczyn, jądro rozpinają te
  $v$, których $\sigma$ jest zerem, najlepsza aproksymacja niskiego rzędu to
  kilka pierwszych składników, a $\kappa$ to pierwsza podzielona przez
  ostatnią.}
\end{summarybox}

\canyou

\begin{testexercises}
\item $A$ jest $5 \times 3$. Podaj, ile ma wartości osobliwych, i powiedz, czy
  może mieć wartości własne.
  \answerto{Trzy \dash{} tyle, ile wynosi mniejszy z wymiarów. Wartości własnych
  nie może mieć w ogóle: $Av$ żyje w przestrzeni pięciowymiarowej, a $v$ w
  trójwymiarowej, więc $Av = \lambda v$ nie jest zdaniem, które da się
  wypowiedzieć.}
\item $A$ ma wartości osobliwe $6$, $3$ i $1$. Podaj $\kappa(A)$,
  $\kappa(A\T A)$ oraz z grubsza, ile cyfr dziesiętnych kosztuje każdy.
  \answerto{$6$ i $36$. Jakieś $\num{0.8}$ cyfry i jakieś $\num{1.6}$ \dash{}
  $\logb{10}6$ i $\logb{10}36$. Chodzi nie o wielkość żadnego z nich, tylko o
  to, że drugi jest dwa razy większy od pierwszego, zawsze.}
\item Obcięcie $A$ do rzędu $2$ zostawia błąd $1$. Co to mówi o $\sigma_3$ i co
  to mówi o każdej innej macierzy rzędu $2$?
  \answerto{$\sigma_3 = 1$, dokładnie, bo zostawiony błąd to to, co wyrzucono. A
  na mocy Eckarta--Younga żadna inna macierz rzędu $2$ nie jest bliżej, więc $1$
  nie jest błędem tej metody, tylko najlepszym dostępnym.}
\item Dwie macierze mają widma $(100, 99, 98, 97)$ oraz $(100, 3, 2, 1)$.
  Powiedz, którą warto aproksymować rzędem $1$ i dlaczego.
  \answerto{Drugą. Rząd $1$ zachowuje $100^{2}$ z
  $100^{2}+3^{2}+2^{2}+1^{2}$, czyli ponad $99\%$; pierwsza zachowuje jakąś
  ćwiartkę, bo jej cztery kierunki liczą się niemal tak samo. \textbf{Warunkiem
  jest zanik, a nie rozmiar.}}
\item Ktoś mówi ci, że solwer zwrócił współczynniki, które \enquote{wyglądają
  rozsądnie}, dla źle uwarunkowanego dopasowania najmniejszych kwadratów.
  Powiedz, dlaczego to nie uspokaja.
  \answerto{Bo porażka uwarunkowania nie produkuje bzdur. Produkuje wiarygodne
  liczby o właściwym znaku i rzędzie wielkości, błędne na cyfrach, których nikt
  nie sprawdza. \enquote{Wygląda rozsądnie} jest dokładnie tym objawem.}
\item Dlaczego macierz z zerową wartością osobliwą nie ma odwrotności, w języku
  tego programu?
  \answerto{Bo $A^{+}$ dzieli przez każdą wartość osobliwą, a dzielenie przez
  zero jest miejscem, w którym to się kończy. Równoważnie: $\Sigma$ spłaszcza
  ten kierunek, więc odwzorowanie nie jest różnowartościowe, co jest jądrem z
  Programu~\ref{prog:P08} i zerowym wyznacznikiem z Programu~\ref{prog:P09}.}
\end{testexercises}

\begin{furtherproblems}
\item Pokaż, że $\kappa(A) \ge 1$ dla każdej macierzy, i powiedz, które macierze
  osiągają $1$.
  \answerto{$\sigma_{\max} \ge \sigma_{\min}$ z definicji, więc stosunek wynosi
  co najmniej $1$. Równość oznacza, że każda wartość osobliwa jest ta sama, więc
  $\Sigma$ jest wielokrotnością identyczności, a $A$ obrotem przeskalowanym
  stałą \dash{} co rozciąga każdy kierunek jednakowo i nie może zniekształcić
  błędu względnego wcale.}
\item Program~\ref{prog:P10} powiedział, że ćwierć obrotu nie ma żadnej
  rzeczywistej wartości własnej. Podaj jej wartości osobliwe.
  \answerto{Obie $1$. $A\T A = I$ dla każdego obrotu, więc każda wartość
  osobliwa wynosi $1$, a $\kappa = 1$. SVD nie ma najmniejszej trudności z
  macierzą, która pokonała wartości własne całkowicie, i to jest teza sekcji 1
  uczyniona konkretną.}
\item Macierz zostaje pomnożona przez $1000$. Powiedz, co dzieje się z jej
  wartościami osobliwymi, z jej wskaźnikiem uwarunkowania i z liczbą cyfr, jakie
  traci rozwiązanie z nią.
  \answerto{Każda wartość osobliwa zostaje pomnożona przez $1000$; wskaźnik
  uwarunkowania nie zmienia się, bo jest stosunkiem; i tracone cyfry też się nie
  zmieniają. To pułapka z wyznacznikiem z Programu~\ref{prog:P10} od drugiej
  strony: skalowanie rusza wyznacznikiem ogromnie, a uwarunkowaniem wcale, i
  dlatego $\kappa$ jest uczciwą miarą \emph{prawie osobliwej}, a $\det$ nie
  jest.}
\item Dlaczego pseudoodwrotność macierzy z maleńką wartością osobliwą daje
  odpowiedź ogromnej długości i jaki jest standardowy środek zaradczy?
  \answerto{Bo dzieli przez tę wartość osobliwą, więc każda składowa prawej
  strony wzdłuż odpowiadającego kierunku zostaje wzmocniona jej odwrotnością.
  Środkiem jest obcięcie \dash{} odrzucenie kłopotliwych wartości osobliwych
  przed odwracaniem \dash{} co ogranicza wzmocnienie kosztem odmowy
  odpowiadania w tych kierunkach.}
\item Program~\ref{prog:P08} dowiódł, że wąskie gardło rzędu $r$ ogranicza rząd
  wszystkiego po nim, i odmówił twierdzenia, że głębokie stosy tracą rząd bez
  niego. Powiedz, jaki pomiar rozstrzygnąłby to drugie i dlaczego ta książka go
  nie wykonała.
  \answerto{Wartości osobliwe reprezentacji na każdej głębokości: jeśli zanikają
  coraz stromiej z głębokością, to rząd spada w tym sensie, o który w tezie
  chodzi. Ta książka pomiaru nie wykonała, bo potrzebuje on wytrenowanego
  modelu, którego nie ma i nie pobiera \dash{} a zbudowanie widma i podanie go
  byłoby wytworem, a nie pomiarem.}
\end{furtherproblems}
